\section{Potential Future Works}
Beyond its efficacy in enhancing time-series forecasting accuracy through multimodality extension (Section~\ref{sec:exp}), \model holds significant potential in advancing time series analysis across a wide spectrum. In the following section, we discuss how \model might transform conventional approaches, facilitate novel methodologies, and may broadly impact the time series analysis domain in future works.
\par\noindent\textbf{Multimodal Time-Series Imputation.} Missing values in time series data, caused by sensor failures, system instability, or privacy concerns, pose a significant challenge in analysis. Conventional time-series imputation (TSI) methods~\cite{che2018recurrent, fortuin2020gp, bansal2021missing, bilovs2023modeling} often overlook valuable information captured in textual formats alongside the numerical data. For instance, incident reports, weather conditions, and special events can provide crucial context for imputing missing data points in traffic time series, but current methods fail to effectively incorporate this information. By constructing missing values using existing toolkits~\cite{du2024tsi}, \model can directly serve as multimodal imputation datasets. \model enables the integration of textual contextual information with numerical time series data, opening new avenues for multimodal time series analysis and enhancing imputation accuracy.

\par\noindent\textbf{Multimodal Time-Series Anomaly Detection.} Detecting anomalies in time series data is crucial for identifying unusual patterns that may indicate faults, fraud, or other significant events~\cite{cook2019anomaly, blazquez2021review, zhang2024self}. However, conventional anomaly detection methods~\cite{munir2018deepant, carreno2020analyzing, xu2021anomaly, shen2020timeseries} are limited to expected pattern deviations in numerical data, overlooking valuable information in textual formats. For instance, news articles, social media posts, and market reports can provide critical context that influences financial market behavior and helps identify anomalies not evident from numerical data alone. \model enables a feasible way to support multimodal time series anomaly detection. For example, in influenza (health) dataset, if "flu outbreak" is mentioned in the text data, we can initially label the aligned numerical data as anomalous.

\par\noindent\textbf{Multimodal Foundation Time-Series Models.} 
The introduction of \model, a comprehensive text-numeric TS dataset, is expected to significantly advance multimodal-based TS methods, including the development of multimodal foundation TS models beyond existing unimodal models~\cite{kamarthi2023large}. \model will facilitate further exploration with more complex and informative prompts, enhancing the performance and capabilities of fine-tuning methods. Additionally, it might spur research and development of multimodal models specifically tailored for TSA, an area with limited exploration compared to other domains like vision and video generation~\cite{singh2022flava,wang2023image,xing2024make,du2024learning}.